% Part: first-order-logic
% Chapter: beyond
% Section: higher-order-logic

\documentclass[../../../include/open-logic-section]{subfiles}

\begin{document}

\olfileid{fol}{byd}{hol}

\olsection{Hogere-ordelogica}

De overgang van eerste- naar tweede-ordelogica stelde ons in staat om binnen
de formele taal over verzamelingen objecten uit het eerste-orde!!{domain} te
spreken. Waarom zouden we daar ophouden? Met derde-ordelogica zouden we
bijvoorbeeld verzamelingen van verzamelingen objecten moeten kunnen
behandelen, of wellicht zelfs verzamelingen die zowel objecten als
verzamelingen objecten bevatten. Vierde-ordelogica laat ons vervolgens over
verzamelingen van zulke objecten spreken. Dit idee kan, zoals te verwachten
valt, willekeurig vaak worden herhaald.

In de praktijk gebruikt de !!{formula}ring van hogere-ordelogica vaak functies
in plaats van relaties. (Op de natuurlijke identificaties na is dit verschil
niet wezenlijk.) Uitgaande van enkele basale ``soorten'' $A$, $B$,
$C$,~\dots (die we nu ``typen'' zullen noemen) kunnen we nieuwe typen vormen
door vast te leggen:
\begin{quote}
Als $\sigma$ en $\tau$ eindige typen zijn, dan is ook $\sigma \to \tau$ een
eindig type.
\end{quote}
Typen kunnen worden opgevat als syntactische ``labels'' die de objecten
classificeren die we in ons !!{domain} willen opnemen; $\sigma \to \tau$
beschrijft de objecten die functies zijn van objecten van type~$\sigma$ naar
objecten van type~$\tau$. We zouden bijvoorbeeld een type $\Omega$ van de
waarheidswaarden ``waar'' en ``onwaar'' kunnen invoeren, en een type $\Nat$ van
de natuurlijke getallen. Objecten van type $\Nat \to \Omega$ kunnen dan worden
opgevat als unaire relaties, of deelverzamelingen van $\Nat$; objecten van type
$\Nat \to \Nat$ zijn functies van natuurlijke getallen naar natuurlijke
getallen; en objecten van type $(\Nat \to \Nat) \to \Nat$ zijn
``functionalen'', dat wil zeggen functies van een hoger type die functies als
invoer nemen en getallen als uitvoer geven.

Net als tweede-ordelogica kan hogere-ordelogica worden opgevat als een soort
meersoortige logica, met een soort voor elk type object dat we willen
beschouwen. Het is gewoonlijk echter duidelijker om de syntaxis van de logica
met hogere typen van onderaf op te bouwen. Zo kunnen we de verzameling eindige
typen als volgt inductief definiëren:
\begin{enumerate}
\item $\Nat$ is een eindig type.
\item Als $\sigma$ en $\tau$ eindige typen zijn, dan is ook $\sigma
  \to \tau$ een eindig type.
\item Als $\sigma$ en $\tau$ eindige typen zijn, dan is ook $\sigma \times
  \tau$ een eindig type.
\end{enumerate}
Intuïtief duidt $\Nat$ het type van de natuurlijke getallen aan, duidt
$\sigma \to \tau$ het type aan van functies van $\sigma$ naar $\tau$, en
duidt $\sigma \times \tau$ het type aan van paren objecten waarvan het ene uit
$\sigma$ en het andere uit $\tau$ komt. Vervolgens kunnen we als volgt
inductief een verzameling termen definiëren:
\begin{enumerate}
\item Voor elk type $\sigma$ is er een voorraad variabelen $x$, $y$,
  $z$, \dots van type $\sigma$
\item $\Obj 0$ is een term van type $\Nat$
\item $\Obj S$ (opvolger) is een term van type $\Nat \to \Nat$
\item Als $s$ een term van type $\sigma$ is en $t$ een term van type
  $\Nat \to (\sigma \to \sigma)$, dan is $\Obj{R}_{st}$ een term van type
  $\Nat \to \sigma$
\item Als $s$ een term van type $\tau \to \sigma$ is en $t$ een
  term van type~$\tau$, dan is $s(t)$ een term van type $\sigma$
\item Als $s$ een term van type~$\sigma$ is en $x$ een variabele van
  type~$\tau$, dan is $\lambd[x][s]$ een term van type $\tau \to \sigma$.
\item Als $s$ een term van type~$\sigma$ is en $t$ een term van
  type~$\tau$, dan is $\tuple{s, t}$ een term van type $\sigma \times
  \tau$.
\item Als $s$ een term van type~$\sigma \times \tau$ is, dan is $p_1(s)$ een
  term van type~$\sigma$ en $p_2(s)$ een term van type~$\tau$.
\end{enumerate}
Intuïtief duidt $\Obj{R}_{st}$ de recursief gedefinieerde functie aan waarvoor
\begin{align*}
\Obj{R}_{st}(0) & = s \\
\Obj{R}_{st}(x+1) & = t(x, R_{st}(x)),
\end{align*}
duidt $\tuple{s, t}$ het paar aan met eerste component~$s$ en tweede
component~$t$, en duiden $p_1(s)$ en~$p_2(s)$ respectievelijk het eerste en
tweede element (de ``projecties'') van~$s$ aan. Ten slotte duidt
$\lambd[x][s]$ de functie~$f$ aan die wordt gedefinieerd door
\[
f(x) = s
\]
voor elke~$x$ van type~$\sigma$; onderdeel (6) verschaft dus een vorm van
comprehensie waarmee functies aan de hand van termen kunnen worden
gedefinieerd. !!^{formula}s worden opgebouwd uit !!{identity}suitspraken
$\eq[s][t]$ tussen termen van hetzelfde type, de gebruikelijke
propositionele connectieven en kwantificatie over hogere typen. Als axioma's
van het systeem kan men vervolgens de basisvergelijkingen nemen die de
hierboven gedefinieerde termen beheersen, samen met de gebruikelijke
logicaregels voor kwantoren en !!{identity}.

Als het systeem van eindige typen wordt uitgebreid met een type $\Omega$ van
waarheidswaarden, moeten ook axioma's worden opgenomen die het gebruik ervan
regelen. Met een geschikte aanpak kunnen complexe !!{formula}s zelfs geheel
worden geëlimineerd en vervangen door termen van type $\Omega$!{} Het
bewijssysteem kan dienovereenkomstig worden aangepast. Het resultaat is in
wezen de door Alonzo Church in de jaren dertig ontwikkelde
\emph{eenvoudige typentheorie}.

Net als bij tweede-ordelogica zijn er verschillende versies van de semantiek
voor hogere typen denkbaar. In de volledige versie lopen variabelen van type
$\sigma \to \tau$ over de verzameling van
\emph{alle} functies van de objecten van type~$\sigma$ naar objecten van
type~$\tau$. Zoals te verwachten is deze semantiek te sterk om een volledig,
effectief !!{derivation} systeem toe te laten. Men kan echter een zwakkere
semantiek beschouwen, waarin !!a{structure} bestaat uit verzamelingen elementen
$T_\tau$ voor elk type $\tau$, samen met passende bewerkingen voor toepassing,
projectie enzovoort. Als de bijzonderheden correct worden uitgewerkt, kunnen
voor de hierboven beschreven soorten !!{derivation} systemen
volledigheidsstellingen worden verkregen.

Logica met hogere typen is aantrekkelijk omdat zij een raamwerk biedt waarin
een aanzienlijk deel van de wiskunde op natuurlijke wijze kan worden
ingebed: uitgaande van $\Nat$ kunnen de reële getallen, continue functies
enzovoort worden gedefinieerd. Het raamwerk is ook bijzonder aantrekkelijk in
de context van de intuïtionistische logica, omdat de typen duidelijke
``constructieve'' interpretaties hebben. Er kunnen zelfs constructieve versies
van de semantiek voor hogere typen worden ontwikkeld (gebaseerd op
intuïtionistische in plaats van klassieke logica) die deze constructieve
interpretaties fraai verduidelijken en in veel opzichten interessanter zijn
dan hun klassieke tegenhangers.

\end{document}
